\subsection{Линейная регрессия}

Пусть наблюдаемая случайная величина \( Z \) зависит от другой случайной величины \( X \) --- значения последней мы можем наблюдать. Обозначим через \( f(x) \) функцию, которая отражает зависимость среднего значения \( Z \) от \( X \):
\[
    E[Z\mid X=x]=f(x)
\]

Функция \( f(t) \) называется \textit{линией регрессии} \( Z \) на \( X \), а уравнение \( t=f(x) \) --- \textit{уравнением регрессии}.

Проведём \( n=20 \) экспериментов и рассмотрим выборку \( \vec{X}=(X_1,\ldots,X_n)^T \), для неё получим значения наблюдаемой величины \( \vec{Z}=(t_1,\ldots,t_n)^T \).

Обозначим через \( \varepsilon_i \) разницу между наблюдаемой в \( i \)-м эксперименте случайной величиной и её математическим ожиданием:
\[
    \varepsilon_i=Z_i-E[Z\mid X=x_i]=Z_i-f(x_i)
\]

По условию лабораторной функция \( f(x) \) --- \textbf{линейная регрессия}, и она имеет вид:
\[
    f(x)=\vec{\omega}^T\cdot\vec{\varphi}(x),
\]

где:
\begin{itemize}
    \item \( \vec{\omega}=\begin{pmatrix}
        \omega_0 & \omega_1 & \omega_2
    \end{pmatrix}^T \) --- вектор неизвестных параметров \textit{(коэффициенты линейной регрессии)};
    \item \( \vec{\varphi}(x)=\begin{pmatrix}
        1 & x & x^2
    \end{pmatrix}^T \) --- вектор факторов регрессии.
\end{itemize}

То есть нашу модель можно записать в матричной форме:
\begin{equation}
    \vec{Z}=\Phi^T\vec{\omega}+\vec{\varepsilon},\label{eq:model}
\end{equation}

где \( \Phi(3\times n) \) --- \textit{матрица плана} вида:
\[
    \Phi=\begin{pmatrix}
        \vec{\varphi}(X_1) & \vec{\varphi}(X_2) & \dots & \vec{\varphi}(X_n)
    \end{pmatrix}
\]